\justifying
\NSFCBodyText

\subsubsection{与本项目相关的研究工作积累}

申请人颜丰为新疆大学副教授、硕士研究生导师，长期从事人工智能、机器学习与深度学习、时间序列预测及多模态信息处理研究。其研究积累可为本项目的三个核心环节提供支撑：一是新能源功率预测和时序建模工作有助于理解光伏、风电等材料中的数据结构、运行曲线、工况与参数属性；二是视觉问答、跨模态检索和多模态表示研究可支撑文本、图像、图表和代码资源的联合表示与关系判定；三是智能知识库与问答系统研发经历可支撑资源元数据、来源证据和质量状态的组织。

在新能源数据智能分析方面，申请人作为通讯作者在 \textit{Energy} 发表了光伏功率预测相关研究《DWT-Former: Fusing Wavelet-Based Multi-Scale Features and Transformer-Based Temporal Representations for Photovoltaic Power Forecasting》（2025, 139283）；作为共同第一作者发表了长时序预测研究《WaveKAN: Adaptive Wavelet Decomposition and Kolmogorov-Arnold Network for Long-Term Time Series Forecasting》（\textit{Neurocomputing}, 2026, 133361）。这些工作形成了对时序数据、多尺度特征、运行曲线和数据分析流程的研究基础，可用于本项目中新能源数据与曲线类资源的专业属性建模，但不将已有预测结果直接等同于教学资源质量结论。

上述积累与本项目第一、二项研究内容的对应关系在于：功率预测任务通常需要明确预测对象、输入变量、采样频率、时间窗口、评价口径和运行背景，任何一个字段缺失都可能使曲线、数据说明与算法案例发生错误关联。申请人可据此将时序材料从“图片或数据文件”进一步表示为包含变量、单位、时间尺度和工况的专业资源片段，并在关系判定中构造参数或时序错配的困难负例。这里使用的是既有研究形成的对象理解和实验组织能力，而非声称已经完成了本项目的课程资源构建或质量验证。

在多模态理解与跨模态关联方面，申请人作为通讯作者开展了跨模态检索、视觉问答和遥感视觉问答研究，包括《Feature Fusion Mamba Hashing via Decoupling for Cross-Modal Retrieval》（\textit{IEEE Signal Processing Letters}, 2026）、《RSHR+: Progressive Question-Conditioned Visual Calibration and Structured State-Space Reasoning for Remote Sensing Visual Question Answering》（\textit{Expert Systems with Applications}, 2026, 132953）和《RSSR: Efficient Knowledge Transfer and Deep Spatial Modeling for Remote Sensing Visual Question Answering》（\textit{Computer Vision and Image Understanding}, 2026, 104753）。此外，申请人在视觉问答方向发表了《Knowledge-Aware Image Understanding with Multi-Level Visual Representation Enhancement for Visual Question Answering》（\textit{Machine Learning}, 2024, 113(6):3789--3805）和《SPCA-Net: A Spatial Position Relationship Co-Attention Network for Visual Question Answering》（\textit{The Visual Computer}, 2022, 38(1):3097--3108）。上述工作为本项目的多模态表示、关系匹配、知识增强和误差分析提供方法基础。

这些多模态研究经验能够支撑项目中“先生成候选、再施加专业约束”的技术路线：视觉或图文模型用于提出可能相关的文本、设备图、曲线和代码片段，知识与属性约束用于排除不满足对象、变量、单位、工况和时序条件的候选，来源位置用于使关系可回溯。申请人已有工作覆盖表示学习、跨模态检索和知识增强等方法环节，因此具备组织基线、实施关系分类与开展消融分析的条件；但专业教学资源的课程对齐、出处证据和四层质量效度仍是本项目需要新开展和验证的内容。

申请人主持的相关项目包括自治区高校基本科研业务费科研项目“融合多源数据大模型驱动的光伏发电功率预测方法研究”（XJEDU2026P005，2026.01--2026.12）、自治区天池英才引进计划（青年博士人才）项目“多模态融合技术研究及其在新疆产业集群中的应用”（2025.02--2028.02，经费50万元）、自治区优秀博士后资助项目“面向通用领域的多模态视觉问答关键技术研究”（2024.07--2025.11，经费3万元）以及横向课题“面向政企数字化应用的智能知识库与问答系统研发”（202604140013，2026.06--2026.10，经费12万元）。这些项目在新能源数据分析、多模态理解和知识组织三个方面与本项目形成方法衔接；正式申报表中的完整资助机构、项目类别、批准号和负责人信息以核验材料为准。

这些项目与本申请在研究对象、数据边界和预期结论上均需严格区分：新能源功率预测项目面向预测方法与数据建模，多模态融合和视觉问答项目面向通用多模态方法，智能知识库项目面向政企数字化应用；本项目面向高校新能源课程资源单元的构建质量与可追溯性，不接入生产系统，不处理学生行为数据。它们提供的是算法、数据组织和工程实施经验，而不是本项目假设已经得到验证的替代证据。

\subsubsection{已取得的研究工作成绩}

申请人已形成新能源时序建模、多模态视觉语言理解、跨模态检索和知识增强表示等相关成果积累。其研究在光伏功率预测、长时间序列预测、跨模态检索、视觉问答及小目标检测等方向形成了从数据表示、特征学习到模型验证的技术基础。相关论文和项目经历表明，申请人具备开展模型设计、实验复现、对照分析和研究生协同组织的能力。

申请人具有人工智能和深度学习相关课程的教学经历，可将算法概念、数据分析过程和实验代码组织为课程材料。项目已确定以面向新能源相关专业本科生的“新能源数据分析与功率预测”示范课程模块开展方法验证，并在模块内设置新能源运行数据与变量认知、功率数据预处理与可视化、功率预测建模与实验、结果评价与工程案例4个知识模块。申请人确认可提供自主编写的讲义、课件、实验代码、图表和案例；这些材料须在入池前逐项登记版本、权利状态和可展示范围，故不将其或拟构建规模、专家评价结果表述为既有成果。

为保证研究基础与拟解决问题逐项对位，申请人的研究积累主要支撑三类工作：时序预测与新能源数据处理支撑曲线、变量与工况属性建模；跨模态检索与视觉问答支撑资源片段表示、候选关系生成和错误分析；知识库研发与课程教学经历支撑来源字段、资源元数据和课程表达的设计。三类基础共同使项目可以从受控资源输入开始开展比较实验，但无法替代课程授权、专家名单和独立盲评等需在立项后落实的条件，相关事项已在工作条件中单列风险与解决途径。

\subsubsection{项目组能力与任务分工}

申请人已组建由约10名研究生构成的科研团队，团队具有人工智能、深度学习、时间序列预测、多模态信息处理和相关应用研究基础。拟参与本项目的研究生包括孔哲、张建宏、杨铄、孟星宇、桑乐吕、中原和于得武。项目实施中，申请人负责总体科学问题凝练、研究方案设计、关键方法把关、资源合规边界和实验质量控制；研究生在申请人指导下参与公开或授权资源整理、知识属性标注、多模态解析、模型训练、对照与消融实验、原型实现和结果复核。个人任务及工作量按立项后的实际到位情况和研究计划分配，不把尚未发生的分工作为既有成果。

\subsubsection{前期基础与本项目的衔接}

前期新能源预测研究提供了对功率曲线、数据变量、工况和时间序列特征的理解，可支撑产业任务与课程知识中专业属性的表示；多模态问答和跨模态检索研究提供文本—视觉表示、关系匹配和知识增强的算法基础，可支撑设备图、运行曲线、代码和文本之间的关联建模；智能知识库研发基础可支撑资源单元、来源证据与质量状态的结构化组织；人工智能和深度学习课程教学经历可支撑资源单元的课程表达和专家评价设计。

上述基础与本项目的关系是方法和对象层面的衔接，而非直接证明本项目假设已经成立。项目仍需通过产业—课程知识消融、专业属性与来源约束消融、独立专家盲评和跨来源留出测试，验证联合约束、跨模态关联和四层质量门控的增量作用。

项目实施中，申请人将负责统一数据与合规边界、知识与属性定义、模型和评估设计以及关键结果复核；研究生成员在指导下完成资源整理、标注、实验执行和可复现记录。每一阶段均保留输入材料、版本、参数、输出和人工修订痕迹，使前期技术积累能够以可审计的方式转化为本项目的研究能力，而不将团队规模等同于成果保证。
